\documentclass[12pt,a4paper]{article}
\usepackage[utf8]{inputenc}
\usepackage[spanish]{babel}
\usepackage{amsmath}
\usepackage{amsfonts}
\usepackage{amssymb}
\usepackage{graphicx}
\usepackage{lmodern}
\usepackage{xcolor}
\usepackage{fancyhdr}
\usepackage[left=2cm,right=2cm,top=3cm,bottom=3cm]{geometry}

\lhead{XXXVIII Coloquio de la Sociedad Matemática Peruana}
\rhead{}

\renewcommand{\headrulewidth}{0.15pt}
\renewcommand{\footrulewidth}{0.15pt}
\begin{document}
\pagestyle{fancy}

\begin{center}
{\large  \textbf{Homología de espacios de intersección y simetría especular}}
\end{center}


\begin{center}
{{Abraham Rojas Vega}\\
{\it {abraham.rojas@usp.br}\\
\small{Universidade de São Paulo}}}\\
\end{center}


\begin{abstract}

La teoría de espacios de intersección (HI), introducida por Markus Banagl \cite{Referencia1}, proporciona un marco novedoso para entender la topología de espacios singulares. Sea $X$ una variedad con singularidades aisladas, $HI_*(X)$ es definida como la homologia de la variedad $IX$, construida reemplazando los links de las singularidades por sus aproximaciones de Moore. Posee propiedades semajantes a la homología de intersección (IH), en particular satisface la dualidad de Poincaré. 

Recientemente, Agustín y Fernández de Bobadilla \cite{Referencia2} dieron una versión axiomática de la HI, similar a la de la IH, mostrando una conexión más evidente entre ambas. En esta charla comentaremos sobre problemas abiertos en esta área, buscando motivar a los asistentes a embarcarse en la investigación de estos temas.

    
\end{abstract}


\noindent\textbf{Palabras Claves: Llenar aquí algunas palabras claves} 



\begin{thebibliography}{10}
\bibitem{Referencia1}
Banagl, M. (2009). Singular spaces and generalized Poincaré complexes. Electronic Research Announcements in Mathematical Sciences, 16(0), 63–73.


\bibitem{Referencia2}
Agustín Vicente, M., \& Fernandez De Bobadilla, J. (2020). Intersection Space Constructible Complexes. Documenta Mathematica, 25, 1653–1725.

\end{thebibliography}
\end{document}